\subsection{Affine charts of the normalized projective model}
\label{subsec:projective-charts}

Let $R=k[X_0,\ldots,X_n]$, with $k$ trivially valued, and fix an index $j$.
Put $A_j=k[U_i:i\ne j]$. The projective model is the space of graded
seminorms on $R$ normalized by $\max_i|X_i|=1$. The construction below
identifies its coordinate opens with affine seminorm spaces. Identifying this
atlas with the analytification of an arbitrary scheme and gluing coherent
sheaves are further geometric constructions.

\begin{definition}[Coordinate open and ratio map {\cite[\S2.1]{FHX}}]
\label{charts:ratio-map}
\lean{EnvelopeProperty.ProjBerk.chartOpen,EnvelopeProperty.ProjBerk.chartRatioHom,
EnvelopeProperty.ProjBerk.coordinateChartMap}
\leanok
\uses{def:pn-an,lem:principal-open}
Let $V_j=\{v:|X_j|_v\ne0\}$ in the normalized projective model.
The homomorphism $\beta_j:A_j\to R[X_j^{-1}]$ sends $U_i$ to
$X_i/X_j$. Extending $v$ to $R[X_j^{-1}]$ and restricting along
$\beta_j$ defines a point $c_j(v)$ of $\operatorname{Berk}_k(A_j)$.
\end{definition}

\begin{lemma}[Coordinate opens and forward chart]
\label{charts:forward-chart}
\lean{EnvelopeProperty.ProjBerk.isOpen_chartOpen,
EnvelopeProperty.ProjBerk.exists_mem_chartOpen,
EnvelopeProperty.ProjBerk.continuous_coordinateChartMap,
EnvelopeProperty.ProjBerk.coordinateChartMap_val_X}
\leanok
\uses{charts:ratio-map,lem:principal-open}
Each $V_j$ is open, the $V_j$ cover the normalized projective model,
and $c_j$ is continuous. For $i\ne j$,
\[
 |U_i|_{c_j(v)}=|X_i|_v/|X_j|_v.
\]
More generally, the value of $a\in A_j$ is its value after substituting
the coordinate ratios in the localized seminorm.
\end{lemma}
\begin{proof}
Evaluation on $X_j$ is continuous and the nonzero locus in
$\R_{\ge0}$ is open. Normalization and finiteness of the coordinates
give an index with $|X_j|=1$, proving that the opens cover.
The principal-localization identification has continuous inverse,
so localization of $v$ is continuous on $V_j$. Pullback along
$\beta_j$ is continuous, proving continuity of $c_j$.
The localization formula evaluates $X_i/X_j$ as the displayed ratio;
the same pullback definition evaluates every polynomial in these ratios.
\end{proof}

\begin{lemma}[Homogeneous substitution]
\label{charts:homogeneous-substitution}
\lean{EnvelopeProperty.ProjBerk.chartHomogenize,
EnvelopeProperty.ProjBerk.coordinateDehom,
EnvelopeProperty.ProjBerk.chartHomogenize_homogeneousComponent,
EnvelopeProperty.ProjBerk.chartHomogenize_coeff,
EnvelopeProperty.ProjBerk.chartRatioHom_dehom_homogeneousComponent_mul_pow}
\leanok
\uses{charts:ratio-map}
Define $H_j:R\to A_j[T]$ by $X_j\mapsto T$ and
$X_i\mapsto U_iT$ for $i\ne j$, and define
$D_j:R\to A_j$ by $X_j\mapsto1$, $X_i\mapsto U_i$.
If $F_d$ is the homogeneous component of degree $d$ of $F$, then
\[
 H_j(F_d)=D_j(F_d)T^d,\qquad
 \operatorname{coeff}_{T^d}H_j(F)=D_j(F_d).
\]
In $R[X_j^{-1}]$,
$\beta_j(D_j(F_d))X_j^d=F_d$.
These identities hold over an arbitrary commutative coefficient ring in
place of $k$.
\end{lemma}
\begin{proof}
For a monomial $a\prod_iX_i^{m_i}$ of total degree $d$,
substitution gives
$a\prod_{i\ne j}U_i^{m_i}T^{\sum_im_i}
=a\prod_{i\ne j}U_i^{m_i}T^d$.
Homogeneous polynomials are sums of precisely these monomials, giving
the first identity. Summing over the finitely many homogeneous
components and taking the coefficient of $T^d$ gives the second identity.
Substitution of $X_i/X_j$
into the dehomogenized monomial yields the original monomial multiplied
by $X_j^{-d}$; multiplication by $X_j^d$ cancels this factor.
Addition gives the identity for $F_d$.
\end{proof}

\begin{lemma}[Normalized inverse construction]
\label{charts:normalized-inverse}
\lean{EnvelopeProperty.ProjBerk.rawChartPoint_isGraded,
EnvelopeProperty.ProjBerk.rawChartPoint_normalized,
EnvelopeProperty.ProjBerk.coordinateChartInv,
EnvelopeProperty.ProjBerk.continuous_coordinateChartInv,
EnvelopeProperty.ProjBerk.rawChartPoint_val_homogeneousComponent}
\leanok
\uses{charts:homogeneous-substitution,lem:gauss-point}
For $y\in\operatorname{Berk}_k(A_j)$, put
\[
 M(y)=\max\{1,\max_{i\ne j}|U_i|_y\},\qquad r(y)=M(y)^{-1}.
\]
Pullback along $H_j$ of the Gauss seminorm of radius $r(y)$ defines
a normalized graded point $b_j(y)\in V_j$, continuously in $y$.
For a homogeneous component $F_d$,
$|F_d|_{b_j(y)}=|D_j(F_d)|_y r(y)^d$.
\end{lemma}
\begin{proof}
Since $M(y)\ge1$, the radius is positive. The coefficient formula
for $H_j(F)$ and the Gauss formula give
$|F|_{b_j(y)}=\max_d|D_j(F_d)|_y r(y)^d$, proving gradedness
and the displayed homogeneous component value. Its coordinate values
are $r(y)$ for $X_j$ and $|U_i|_yr(y)$ for $i\ne j$. Their maximum
is $M(y)r(y)=1$, proving normalization and membership in $V_j$.
The finite maximum defining $M$ is continuous and positive, so its
inverse $r$ is continuous. Joint continuity of the Gauss seminorm and
continuous pullback prove continuity of the constructed point.
\end{proof}

\begin{theorem}[Inverse chart and homeomorphism]
\label{charts:homeomorphism}
\lean{EnvelopeProperty.ProjBerk.coordinateChartHomeomorph,
EnvelopeProperty.ProjBerk.coordinateChartMap_coordinateChartInv,
EnvelopeProperty.ProjBerk.coordinateChartInv_coordinateChartMap}
\leanok
\uses{charts:forward-chart,charts:homogeneous-substitution,charts:normalized-inverse,
charts:localized-homogeneous-substitution,charts:localized-ratio-cancellation}
The maps $c_j$ and $b_j$ are mutually inverse homeomorphisms.
\end{theorem}
\begin{proof}
Its ratios have the original affine seminorm: localization sends $X_j$
to an invertible element of value $r(y)$, and substituting $X_i/X_j$
into $H_j$ cancels $T$, giving $U_i$. Hence $c_j(b_j(y))=y$ on
every element of $A_j$.
Conversely let $v\in V_j$. The ratio formula and normalization give
$M(c_j(v))=|X_j|_v^{-1}$ and $r(c_j(v))=|X_j|_v$.
The localization identity in the preceding lemma gives
$|D_j(F_d)|_{c_j(v)}|X_j|_v^d=|F_d|_v$.
The homogeneous evaluation formula therefore shows that
$b_j(c_j(v))$ and $v$ agree on every homogeneous form. Both seminorms
are graded, so their values on every polynomial are the maxima of their
homogeneous component values and hence agree. This proves the other
inverse identity. The two maps are continuous by the forward chart
and normalized inverse construction lemmas.
\end{proof}

